\section{Preliminaries: On-Policy Distillation}
\label{sec:prelim}
\textbf{Multi-turn agent interaction.}
We consider a student agent that interacts with an external environment over
multiple turns. At turn $t$, the agent observes $o_t$, conditions on the full
history
\[
h_t=(x,o_1,r_1,o_2,r_2,\ldots,o_t),
\]
and generates a model response $r_t\sim\pi_\theta(\cdot\mid h_t)$. The response
may contain reasoning text and tool arguments. The
environment then maps the executable part of $r_t$ to the next observation
$o_{t+1}$. A rollout is therefore an interaction trace
\[
\tau=(x,o_1,r_1,o_2,r_2,\ldots,o_T,r_T),
\]
which terminates when the task is completed, the environment returns a terminal
state, or a maximum horizon is reached. This turn structure matters because an
early student action changes later observations and thus changes the future
contexts on which the teacher will be queried.

\textbf{Multi-turn OPD.}
We use on-policy distillation (OPD) \citep{lu2025onpolicydistillation}: the student
samples its own rollout $\tau\sim\pi_\theta(\cdot\mid x)$ for a prompt
$x\sim p_{\mathrm{data}}$, and a frozen teacher $\pi_T$ provides token-level
supervision on the student-visited prefixes. For token position $i$ in the
concatenated model responses, let $s_i$ denote the full prefix before that token,
including previous observations, previous student responses, and the current
partial response. This differs from offline distillation because the teacher is
queried on histories induced by the current student rather than only on
teacher-generated demonstrations.

For a response-token mask $m_i$, we use reverse KL as the OPD objective:
\begin{equation}
\mathcal{L}_{\mathrm{OPD}}(\theta)
=
\mathbb{E}_{x,\tau}
\left[
\frac{1}{\sum_i m_i}
\sum_{i=1}^{L}
m_i\,
\Dkl\!\left(
\pi_\theta(\cdot\mid s_i)\,\|\,\pi_T(\cdot\mid s_i)
\right)
\right].
\label{eq:opd-rkl}
\end{equation}
